\section{Conservación variacional y acción orientada}
\label{nuclear:7}
\subsection{Simetría de la acción elíptica y carga de Noether}

Se fija una sección positiva de acción \(\hbar\), el par angular \(A,C^*\) en una misma unidad, \(A\ne0\), \(|C^*/A|<1\), y

\[
\eta=\operatorname{artanh}(C^*/A).
\]

El plano realizado \(\mathcal P_\eta\) tiene coordenadas \((Q,P)\), ambas de dimensión raíz de acción, con

\[
\lambda=P\,dQ,\qquad \Omega=dQ\wedge dP=-d\lambda.
\]

La construcción de \cref{iv:ellipse:section} da

\begin{equation}
\mathcal I_\eta(Q,P)=\frac12(e^{-\eta}Q^2+e^\eta P^2),
\qquad H_\eta=\omega\mathcal I_\eta,\qquad \omega>0.
\label{nuc07:eq:1}
\end{equation}

La elipse HMT es \(\mathcal I_\eta=\hbar\). Su área es \(2\pi\hbar\); \(\mathcal I_\eta\) tiene dimensión de acción y \(H_\eta\) de energía. Definimos las matrices ya realizadas

\begin{equation}
M_\eta=\operatorname{diag}(e^{\eta/2},e^{-\eta/2}),\quad
J_2=\begin{pmatrix}0&-1\\1&0\end{pmatrix},\quad
R_\eta(\varepsilon)=M_\eta e^{-\varepsilon J_2}M_\eta^{-1}.
\label{nuc07:eq:2}
\end{equation}

\begin{proposition}
El grupo \(R_\eta\) es una simetría canónica de la acción

\begin{equation}
\mathscr S_\eta[z]=\int_{t_a}^{t_b}
\left(P\dot Q-\omega\mathcal I_\eta(Q,P)\right)dt
\label{nuc07:eq:3}
\end{equation}

y su carga de Noether es \(\mathcal I_\eta\).
\end{proposition}

\begin{proof}
El determinante de \(M_\eta\) es uno y \(M_\eta^{\mathsf T}J_2M_\eta=J_2\). Además,
\(\mathcal I_\eta(M_\eta y)=\|y\|^2/2\), invariante bajo \(e^{-\varepsilon J_2}\).
El generador

\[
X_\eta=e^\eta P\,\partial_Q-e^{-\eta}Q\,\partial_P
\]

satisface \(\iota_{X_\eta}\Omega=d\mathcal I_\eta\). Por la fórmula de Cartan,

\[
\mathcal L_{X_\eta}\lambda=dF_\eta,\qquad
F_\eta=\lambda(X_\eta)-\mathcal I_\eta=e^\eta P^2-\mathcal I_\eta.
\]

Como \(X_\eta H_\eta=0\), la variación de la densidad de acción es \(dF_\eta/dt\). Su carga es

\[
\frac{\partial L}{\partial\dot Q}X_\eta Q-F_\eta
=Pe^\eta P-F_\eta=\mathcal I_\eta.
\]

Las ecuaciones \(\dot Q=\omega e^\eta P\), \(\dot P=-\omega e^{-\eta}Q\) verifican directamente \(d\mathcal I_\eta/dt=0\).
\end{proof}

La parametrización positiva de la elipse es \(Q=a_S\cos\theta\), \(P=b_S\sin\theta\). Estas ecuaciones le asignan \(\dot\theta=-\omega\). Por tanto, en una vuelta recorrida según el flujo dinámico, \(\oint P\,dQ=2\pi\mathcal I_\eta\); con \(\theta\) creciente, la integral es \(-2\pi\mathcal I_\eta\). Es la distinción de orientación establecida en \cref{iv:ellipse:area}. En el contorno radional, el módulo es \(h=2\pi\hbar\).

\subsection{Cambio de forma: transporte y conexión determinados}

Para dos formas ya publicadas \(\eta,\eta'\) y un avance de fase \(\vartheta\), se define

\begin{equation}
\mathcal T_{\eta',\eta}^{\vartheta}:\mathcal P_\eta\to\mathcal P_{\eta'},
\qquad
\mathcal T_{\eta',\eta}^{\vartheta}
=M_{\eta'}e^{-\vartheta J_2}M_\eta^{-1}.
\label{nuc07:eq:4}
\end{equation}

\begin{proposition}
Este transporte es simpléctico y cumple

\begin{equation}
\mathcal I_{\eta'}(\mathcal T_{\eta',\eta}^{\vartheta}z)
=\mathcal I_\eta(z),\qquad
\mathcal T_{\eta'',\eta'}^{\vartheta_2}
\mathcal T_{\eta',\eta}^{\vartheta_1}
=\mathcal T_{\eta'',\eta}^{\vartheta_1+\vartheta_2}.
\label{nuc07:eq:5}
\end{equation}
\end{proposition}

\begin{proof}
Todos los factores son simplécticos. La primera igualdad se reduce a la invariancia de \(\|y\|^2/2\) bajo una rotación. En la segunda se cancelan las matrices contiguas \(M_{\eta'}^{-1}M_{\eta'}\) y se suman las fases.
\end{proof}

La fórmula recibe las formas y los avances del lector; no selecciona sus valores. La lista completa de transiciones puede retenerse junto con el transporte compuesto: dos listas que den el mismo mapa compuesto no se identifican como historias.

\begin{proposition}
Una realización diferenciable \(\eta(t),\vartheta(t)\in C^1\), con \(\dot\vartheta=\omega(t)\), del transporte \eqref{nuc07:eq:4} tiene generador

\begin{equation}
\boxed{
H_{\rm tr}(Q,P,t)=
\omega(t)\mathcal I_{\eta(t)}(Q,P)
+\frac{\dot\eta(t)}2 QP.}
\label{nuc07:eq:6}
\end{equation}

Este generador conserva \(\mathcal I_{\eta(t)}(z(t))\).
\end{proposition}

\begin{proof}
Al derivar
\(z(t)=M_{\eta(t)}e^{-\vartheta(t)J_2}y_0\) se obtiene

\[
\dot Q=\omega e^\eta P+\frac{\dot\eta}2Q,\qquad
\dot P=-\omega e^{-\eta}Q-\frac{\dot\eta}2P.
\]

Son las ecuaciones de Hamilton de \eqref{nuc07:eq:6}. El término de conexión realiza exactamente \(\dot M_\eta M_\eta^{-1}z\), y su Hamiltoniano queda determinado salvo una función exclusiva de \(t\). En efecto,

\begin{equation}
\begin{aligned}
\frac{d\mathcal I_\eta}{dt}
&=\frac{\dot\eta}{2}(-e^{-\eta}Q^2+e^\eta P^2)
+\left\{\mathcal I_\eta,\frac{\dot\eta}{2}QP\right\}\\
&=\frac{\dot\eta}{2}(-e^{-\eta}Q^2+e^\eta P^2)
+\frac{\dot\eta}{2}(e^{-\eta}Q^2-e^\eta P^2)=0.
\end{aligned}
\label{nuc07:eq:7}
\end{equation}

Equivalentemente, \(Q=e^{\eta/2}q\), \(P=e^{-\eta/2}p\) dan

\begin{equation}
P\dot Q-H_{\rm tr}
=p\dot q-\frac{\omega(t)}2(q^2+p^2).
\label{nuc07:eq:8}
\end{equation}

La simetría de rotación y su carga se transportan exactamente.
\end{proof}

Aquí se ha derivado el término de conexión del transporte, no de una fuerza externa. El Hamiltoniano LC de forma fija conserva \(\mathcal I_\eta\). Durante una deformación, omitir el segundo término de \eqref{nuc07:eq:6} deja sin cancelar el primer término de \eqref{nuc07:eq:7}. La conservación variacional exige transportar también la forma canónica.

La proposición determina el generador de la interpolación declarada; no asigna una interpolación física única a toda historia TPK. La identidad discreta \eqref{nuc07:eq:4}–\eqref{nuc07:eq:5} funciona sin interpolación.

\subsection{Orientación y frontera de la acción}

Las transformaciones de \cref{iv:ellipse:reciprocidad} distinguen \(S(Q,P)=(P,Q)\) y \(J_2(Q,P)=(-P,Q)\). Ambos intercambian la forma \(\eta\leftrightarrow-\eta\), pero

\begin{equation}
S^*\Omega=-\Omega,\quad J_2^*\Omega=\Omega,\qquad
S^*\lambda=-\lambda+d(PQ),\quad
J_2^*\lambda=\lambda-d(PQ).
\label{nuc07:eq:9}
\end{equation}

Asimismo,

\begin{equation}
S R_\eta(\varepsilon)S^{-1}=R_{-\eta}(-\varepsilon),\qquad
J_2R_\eta(\varepsilon)J_2^{-1}=R_{-\eta}(\varepsilon).
\label{nuc07:eq:10}
\end{equation}

Se comprueban usando \(SM_\eta=M_{-\eta}S\), \(J_2M_\eta=M_{-\eta}J_2\), \(SJ_2S^{-1}=-J_2\). Para una curva abierta orientada \(z:[a,b]\to\mathcal P_\eta\),

\begin{equation}
\int_{J_2z}\lambda=\int_z\lambda-[PQ]_a^b,\qquad
\int_{Sz}\lambda=-\int_z\lambda+[PQ]_a^b.
\label{nuc07:eq:11}
\end{equation}

Los términos de frontera se anulan en un ciclo y se cancelan al concatenar extremos compatibles. La energía cuadrática por sí sola no determina estos signos ni esos términos.

La monodromía nonádica registra retorno de fase y \(q_{\rm mem}\mapsto q_{\rm mem}+1\). La constancia de una carga no obliga al contador a permanecer fijo. En la realización elíptica, una vuelta conserva \(\mathcal I_\eta\) mientras incrementa la integral acumulada orientada en \(2\pi\mathcal I_\eta\). Carga, acción integrada, fase y registro tienen tipos distintos. No se identifica aquí una vuelta angular con cualquier ciclo TPK sin especificar su lector.

\subsection{Carga variacional del reloj y balance entre lectura y memoria}

\subsubsection{Una carga de acción explícita del reloj}

El calendario se realiza sobre \(\mathcal H_{\rm clk}=\mathbb C^{108}\), con proyectores Fourier \(P_k=|\widetilde k\rangle\langle\widetilde k|\), y determina

\begin{equation}
N_{108}=\sum_{k=0}^{107}kP_k,\quad
H_{\rm clk}=\hbar\omega_0N_{108},\quad
\omega_0=\frac{2\pi}{108t_0}.
\label{nuc07:eq:12}
\end{equation}

Su propagador durante \(t_0\) es el desplazamiento de una celda. La acción del reloj es

\begin{equation}
\mathscr S[\psi]=\int
\left[\frac{i\hbar}{2}
(\langle\psi,\dot\psi\rangle-\langle\dot\psi,\psi\rangle)
-\langle\psi,H_{\rm clk}\psi\rangle\right]dt.
\label{nuc07:eq:13}
\end{equation}

El producto es antilineal en la primera variable. La ecuación es \(i\hbar\dot\psi=H_{\rm clk}\psi\).

\begin{proposition}
Para \(B=B^*\) con \([B,H_{\rm clk}]=0\), la simetría \(\psi\mapsto e^{-i\varepsilon B}\psi\) tiene carga

\begin{equation}
\mathcal Q_B(\psi)=\hbar\langle\psi,B\psi\rangle.
\label{nuc07:eq:14}
\end{equation}
\end{proposition}

\begin{proof}
Al localizar el parámetro, \(\delta\psi=-i\varepsilon(t)B\psi\), se obtiene
\(\delta L=\hbar\dot\varepsilon\langle\psi,B\psi\rangle\).
La integración por partes en las ecuaciones variacionales da \(d\mathcal Q_B/dt=0\). Directamente,

\[
\frac{d}{dt}\langle\psi,B\psi\rangle
=\frac{i}{\hbar}\langle\psi,[H_{\rm clk},B]\psi\rangle=0.
\]

Con \(B=I\) se conserva norma; con \(B=N_{108}\) la carga es la acción \(E_{\rm clk}/\omega_0\).
\end{proof}

La forma simpléctica del espacio real subyacente es
\(\Omega_{\rm clk}(u,v)=2\hbar\operatorname{Im}\langle u,v\rangle\);
su uno-forma es
\(\Theta_\psi(v)=-\hbar\operatorname{Im}\langle\psi,v\rangle\),
con \(\Omega_{\rm clk}=-d\Theta\).

\subsubsection{Balance exacto de carga}

Sean \(\mathcal H,\mathcal K\) espacios hermíticos finitos, \(J:\mathcal H\to\mathcal K\) isometría, \(U\) unitario sobre \(\mathcal K\), y

\begin{equation}
\Pi=JJ^*,\quad T=J^*UJ,\quad
\eta_{\rm mem}=(I-\Pi)UJ.
\label{nuc07:eq:15}
\end{equation}

El operador \(\eta_{\rm mem}\) es distinto del parámetro escalar \(\eta\).

\begin{proposition}
Si

\begin{equation}
\widetilde B=\widetilde B^*,\quad U^*\widetilde BU=\widetilde B,
\quad [\widetilde B,\Pi]=0,\quad B=J^*\widetilde BJ,
\label{nuc07:eq:16}
\end{equation}

entonces

\begin{equation}
B=T^*BT+\eta_{\rm mem}^*\widetilde B\eta_{\rm mem},
\qquad
\mathcal Q_B(\psi)=
\mathcal Q_B(T\psi)+\mathcal Q_{\widetilde B}(\eta_{\rm mem}\psi).
\label{nuc07:eq:17}
\end{equation}
\end{proposition}

\begin{proof}
Se expande \(UJ\psi=JT\psi+\eta_{\rm mem}\psi\) en
\(J^*U^*\widetilde BUJ=J^*\widetilde BJ\).
La conmutación con \(\Pi\) anula los términos cruzados. Evaluar la identidad operatoria y multiplicar por \(\hbar\) da el balance de acción.
\end{proof}

Por iteración,

\begin{equation}
\mathcal Q_B(\psi)=\mathcal Q_B(T^n\psi)
+\sum_{j=0}^{n-1}\mathcal Q_{\widetilde B}
(\eta_{\rm mem}T^j\psi).
\label{nuc07:eq:18}
\end{equation}

Esta suma registra complementos sucesivos de la dinámica comprimida. La dinámica completa tiene otra fórmula exacta, con \(z_n=U^nJ\psi\):

\begin{equation}
\mathcal Q_B(\psi)=
\mathcal Q_B(J^*z_n)+
\mathcal Q_{\widetilde B}((I-\Pi)z_n).
\label{nuc07:eq:19}
\end{equation}

No se identifica \(T^n\) con \(J^*U^nJ\) sin otra propiedad: el estado completo puede reabsorber memoria. Si \(B,\widetilde B\ge0\), todos los sumandos de \eqref{nuc07:eq:18} son no negativos; para cargas firmadas permanece la igualdad, sin esa interpretación de positividad.

Si \(U\) es propagador de una acción como \eqref{nuc07:eq:13}, la conmutación de \(\widetilde B\) con su Hamiltoniano aporta la conservación de Noether. La condición de reducción por \(\Pi\) especifica las lecturas que separan carga sin términos de interferencia.

Sin imponer \eqref{nuc07:eq:16}, la identidad general conserva los términos adicionales:

\begin{equation}
\begin{aligned}
&\mathcal Q_B(T\psi)+\mathcal Q_{\widetilde B}(\eta_{\rm mem}\psi)
+2\hbar\operatorname{Re}
\langle JT\psi,\widetilde B\eta_{\rm mem}\psi\rangle\\
&=\mathcal Q_B(\psi)+
\hbar\langle\psi,J^*(U^*\widetilde BU-\widetilde B)J\psi\rangle.
\end{aligned}
\label{nuc07:eq:20}
\end{equation}

La carga física utiliza su lector, por ejemplo \(N_{108}\); una identidad de normas corresponde a \(B=I\). Las secciones \(\hbar_{\rm pre}\) y \(\hbar_{\rm ret}\) conservan el lector de retorno de \cref{iv:action:secciones,iv:action:razon}; la diferencia entre ellas y el complemento de memoria poseen sus respectivos dominios de definición.

\subsection{Balance covariante del tensor total}

La eliminación variacional de la contorsión y la identidad de Noether de la acción covariante, desarrolladas en la proposición~\ref{vii:prop:fuente-metrica-balance}, proporcionan

\begin{equation}
\mathcal T_{\mu\nu}
=T^{\rm phys,LC}_{\mu\nu}+U^{\rm phys,spin}_{\mu\nu},
\qquad
\mathring\nabla^\mu\mathcal T_{\mu\nu}=0.
\label{nuc07:eq:21}
\end{equation}

Su dominio es la carta HMT–MD realizada, con los acoplamientos mantenidos constantes durante la variación. Para un campo de Killing \(\xi\),

\begin{equation}
j^\mu_\xi=\mathcal T^{\mu\nu}\xi_\nu,\qquad
\mathring\nabla_\mu j^\mu_\xi=0.
\label{nuc07:eq:22}
\end{equation}

La prueba usa \eqref{nuc07:eq:21}, la simetría de \(\mathcal T\) y la antisimetría de \(\mathring\nabla_\mu\xi_\nu\). Al integrar, la diferencia entre cargas de las hipersuperficies es el flujo por la frontera lateral.

La conservación pertenece al tensor total; cada componente puede intercambiar carga. El teorema~\ref{vii:thm:intercambio-traza-residual} determina explícitamente \(\mathcal J^{\rm tr}_\nu=\partial_\nu\operatorname{tr}_g(T^{\rm res})/4\) y los balances opuestos de traza y parte de traza nula. Esto se recupera como instancia covariante del criterio de balance conjunto, no como nueva ecuación gravitatoria.
